
        \documentclass[fleqn]{article}      
        \usepackage{amsmath}
        \usepackage{fontspec}
        \usepackage{kotex} % 한국어 지원  

        \begin{document}      
        % 쉬운 문제
\noindent 1. 다음 행렬의 덧셈 결과로 올바른 것은?  
\[
A = \begin{pmatrix} 1 & 2 \\ 3 & 4 \end{pmatrix}, \quad B = \begin{pmatrix} 5 & 6 \\ 7 & 8 \end{pmatrix}
\]  
\begin{itemize}
    \item [(A)] \(\begin{pmatrix} 6 & 8 \\ 10 & 12 \end{pmatrix}\)
    \item [(B)] \(\begin{pmatrix} 5 & 8 \\ 10 & 12 \end{pmatrix}\)
    \item [(C)] \(\begin{pmatrix} 6 & 7 \\ 10 & 12 \end{pmatrix}\)
    \item [(D)] \(\begin{pmatrix} 1 & 2 \\ 3 & 4 \end{pmatrix}\)
\end{itemize}\\\\\\\\

\noindent 2. 다음 행렬의 곱셈 결과로 올바른 것은?  
\[
A = \begin{pmatrix} 1 & 2 \\ 3 & 4 \end{pmatrix}, \quad B = \begin{pmatrix} 5 \\ 6 \end{pmatrix}
\]  
\begin{itemize}
    \item [(A)] \(\begin{pmatrix} 17 \\ 39 \end{pmatrix}\)
    \item [(B)] \(\begin{pmatrix} 12 \\ 34 \end{pmatrix}\)
    \item [(C)] \(\begin{pmatrix} 13 \\ 27 \end{pmatrix}\)
    \item [(D)] \(\begin{pmatrix} 11 \\ 32 \end{pmatrix}\)
\end{itemize}\\\\\\\\

\noindent 3. 다음 행렬의 전치 행렬은 무엇인가?  
\[
A = \begin{pmatrix} 1 & 2 \\ 3 & 4 \end{pmatrix}
\]  
\begin{itemize}
    \item [(A)] \(\begin{pmatrix} 1 & 3 \\ 2 & 4 \end{pmatrix}\)
    \item [(B)] \(\begin{pmatrix} 2 & 4 \\ 1 & 3 \end{pmatrix}\)
    \item [(C)] \(\begin{pmatrix} 1 & 2 \\ 4 & 3 \end{pmatrix}\)
    \item [(D)] \(\begin{pmatrix} 3 & 4 \\ 1 & 2 \end{pmatrix}\)
\end{itemize}\\\\\\\\

\noindent 4. 다음 행렬 \(A\)의 행렬식은 얼마인가?  
\[
A = \begin{pmatrix} 2 & 3 \\ 1 & 4 \end{pmatrix}
\]\\\\\\\\

\noindent 5. 다음 중 \(2 \times 2\) 행렬의 형태로 올바른 것은?  
\[
\begin{pmatrix} a & b \\ c & d \end{pmatrix}
\]  
\begin{itemize}
    \item [(A)] \(a + b + c + d\)
    \item [(B)] \((a, b, c, d)\)
    \item [(C)] \(\begin{pmatrix} a & b \\ c & d \end{pmatrix}\)
    \item [(D)] \(\{a, b, c, d\}\)
\end{itemize}\\\\\\\\

% 중간 문제
\noindent 6. 다음 두 행렬의 곱 \(AB\)는 얼마인가?  
\[
A = \begin{pmatrix} 1 & 0 & 2 \\ 0 & 3 & 1 \end{pmatrix}, \quad B = \begin{pmatrix} 3 & 1 \\ 1 & 2 \\ 0 & 1 \end{pmatrix}
\]\\\\\\\\

\noindent 7. 행렬 \(A\)와 \(B\)의 스칼라 곱 \(2A\)는 얼마인가?  
\[
A = \begin{pmatrix} 4 & -2 \\ 1 & 3 \end{pmatrix}
\]\\\\\\\\

\noindent 8. 다음을 풀어라: \(A\)의 역행렬 구하기  
\[
A = \begin{pmatrix} 1 & 2 \\ 3 & 5 \end{pmatrix}
\]\\\\\\\\

\noindent 9. 행렬 \(A\)의 고유값을 구하라:  
\[
A = \begin{pmatrix} 2 & 1 \\ 1 & 2 \end{pmatrix}
\]\\\\\\\\

\noindent 10. 벡터 \(\mathbf{v}\)의 결과는 얼마인가?  
\[
\mathbf{v} = A \begin{pmatrix} 2 \\ 3 \end{pmatrix}, \quad A = \begin{pmatrix} 4 & 1 \\ 0 & 2 \end{pmatrix}
\]\\\\\\\\

% 어려운 문제
\noindent 11. 다음 조건을 만족하는 행렬 \(A\)를 구하라:  
\[
A^2 - A + I = 0 \quad \text{(단, \(I\)는 항등 행렬)}
\]\\\\\\\\

\noindent 12. 행렬 \(A\)와 \(B\)의 곱과 덧셈의 관계를 보여라:  
\[
A = \begin{pmatrix} 1 & 2 \\ 3 & 4 \end{pmatrix}, \quad B = \begin{pmatrix} 5 & 6 \\ 7 & 8 \end{pmatrix}
\]\\\\\\\\

\noindent 13. 세 개의 행렬 \(A\), \(B\), \(C\)를 사용할 때 \(A + B + C\)를 구하라:  
\[
A = \begin{pmatrix} 1 & 1 \\ 1 & 1 \end{pmatrix}, \quad B = \begin{pmatrix} 2 & 2 \\ 2 & 2 \end{pmatrix}, \quad C = \begin{pmatrix} 3 & 3 \\ 3 & 3 \end{pmatrix}
\]\\\\\\\\

\noindent 14. 다음 행렬의 고유벡터를 구하라:  
\[
A = \begin{pmatrix} 0 & 1 \\ -2 & -3 \end{pmatrix}
\]\\\\\\\\

\noindent 15. 다음 행렬의 계수를 구하라:  
\[
A = \begin{pmatrix} 1 & 2 & 3 \\ 4 & 5 & 6 \\ 7 & 8 & 9 \end{pmatrix}
\]\\\\\\\\

% 답안
1. (A) \\\\\\\\
2. (A) \\\\\\\\
3. (A) \\\\\\\\
4. 2 \\\\\\\\
5. (C) \\\\\\\\
6. \(\begin{pmatrix} 3 & 7 \\ 9 & 6 \end{pmatrix}\) \\\\\\\\
7. \(\begin{pmatrix} 8 & -4 \\ 2 & 6 \end{pmatrix}\) \\\\\\\\
8. \(\begin{pmatrix} 5 & -2 \\ -3 & 1 \end{pmatrix}\) \\\\\\\\
9. \(3, -1\) \\\\\\\\
10. \(\begin{pmatrix} 14 \\ 6 \end{pmatrix}\) \\\\\\\\
11. \(\begin{pmatrix} 1 & 0 \\ 0 & 1 \end{pmatrix}\) \\\\\\\\
12. \(A + B = \begin{pmatrix} 6 & 8 \\ 10 & 12 \end{pmatrix}\) \\\\\\\\
13. \(A + B + C = \begin{pmatrix} 6 & 6 \\ 6 & 6 \end{pmatrix}\) \\\\\\\\
14. \(\begin{pmatrix} 1 \\ 2 \end{pmatrix}\) \\\\\\\\
15. \(0\) \\\\\\\\
      
        \end{document} 
    